
\documentclass[12pt, a4paper, openany, twoside]{book}
\usepackage[left=1in,right=1in,bottom=0.5in,top=0.8in]{geometry}
\usepackage{amsfonts}
\usepackage{amsmath, amssymb}
\usepackage{graphicx}
\usepackage[font=small,labelfont=bf]{caption}
\usepackage{epstopdf}
\usepackage[pdfpagelabels,hyperindex]{hyperref}
\usepackage{xcolor}
\usepackage{amsthm}
\usepackage{float}
\usepackage{pgfplots}
\usepackage{listings}
\usepackage{longtable}
\usepackage{mathrsfs}
\usepackage{mathtools}
\usepackage{hyperref}
\usepackage[most]{tcolorbox}
\usepackage{amssymb}
\usepackage{algorithm}
\usepackage[noend]{algpseudocode}
\usepackage{tikz-cd}
\usepackage{comment}
\usepackage{mathpartir}

\linespread{1.3} 

\hypersetup{
pdftitle={.},
pdfauthor={Hojune Lee},
}

\newcommand{\overbar}[1]{\mkern 1.5mu\overline{\mkern-1.5mu#1\mkern-1.5mu}\mkern 1.5mu}
\DeclareRobustCommand{\stirling}{\genfrac\{\}{0pt}{}}
\newtheorem{thm}{Theorem}[chapter]
\newtheorem{cor}[thm]{Corollary}
\newtheorem{prop}[thm]{Proposition}
\newtheorem{lem}[thm]{Lemma}
\newtheorem{conj}[thm]{Conjecture}
\newtheorem{quest}[thm]{Question}
\newtheorem{ppty}[thm]{Property}
\newtheorem{ppties}[thm]{Properties}
\newtheorem{axiom}[thm]{Axiom}
\newtheorem{claim}[thm]{Claim}
\newtheorem{prob}[thm]{Problem}


\theoremstyle{definition}
\newtheorem{defn}[thm]{Definition}
\newtheorem{defns}[thm]{Definitions}
\newtheorem{con}[thm]{Construction}
\newtheorem{exmp}[thm]{Example}
\newtheorem{exmps}[thm]{Examples}
\newtheorem{notn}[thm]{Notation}
\newtheorem{notns}[thm]{Notations}
\newtheorem{addm}[thm]{Addendum}
\newtheorem{exer}[thm]{Exercise}
\newtheorem{limit}[thm]{Limitation}


\theoremstyle{remark}
\newtheorem{rem}[thm]{Remark}
\newtheorem{rems}[thm]{Remarks}
\newtheorem{warn}[thm]{Warning}
\newtheorem{sch}[thm]{Scholium}
\newenvironment{thmbox}[1][]{\begin{tcolorbox}[colback=yellow!10!white,colframe=red!75!black,title={#1}]}{\end{tcolorbox}}


% newcommand bb
    \newcommand{\BA}{{\mathbb {A}}} \newcommand{\BB}{{\mathbb {B}}}
    \newcommand{\BC}{{\mathbb {C}}} \newcommand{\BD}{{\mathbb {D}}}
    \newcommand{\BE}{{\mathbb {E}}} \newcommand{\BF}{{\mathbb {F}}}
    \newcommand{\BG}{{\mathbb {G}}} \newcommand{\BH}{{\mathbb {H}}}
    \newcommand{\BI}{{\mathbb {I}}} \newcommand{\BJ}{{\mathbb {J}}}
    \newcommand{\BK}{{\mathbb {U}}} \newcommand{\BL}{{\mathbb {L}}}
    \newcommand{\BM}{{\mathbb {M}}} \newcommand{\BN}{{\mathbb {N}}}
    \newcommand{\BO}{{\mathbb {O}}} \newcommand{\BP}{{\mathbb {P}}}
    \newcommand{\BQ}{{\mathbb {Q}}} \newcommand{\BR}{{\mathbb {R}}}
    \newcommand{\BS}{{\mathbb {S}}} \newcommand{\BT}{{\mathbb {T}}}
    \newcommand{\BU}{{\mathbb {U}}} \newcommand{\BV}{{\mathbb {V}}}
    \newcommand{\BW}{{\mathbb {W}}} \newcommand{\BX}{{\mathbb {X}}}
    \newcommand{\BY}{{\mathbb {Y}}} \newcommand{\BZ}{{\mathbb {Z}}}

% newcommand  scr
    \newcommand{\sA}{{\mathscr {A}}} \newcommand{\sB}{{\mathscr {B}}}
    \newcommand{\sC}{{\mathscr {C}}} \newcommand{\sD}{{\mathscr {D}}}
    \newcommand{\sE}{{\mathscr {E}}} \newcommand{\sF}{{\mathscr {F}}}
    \newcommand{\sG}{{\mathscr {G}}} \newcommand{\sH}{{\mathscr {H}}}
    \newcommand{\sI}{{\mathscr {I}}} \newcommand{\sJ}{{\mathscr {J}}}
    \newcommand{\sK}{{\mathscr {K}}} \newcommand{\sL}{{\mathscr {L}}}
    \newcommand{\sN}{{\mathscr {N}}} \newcommand{\sM}{{\mathscr {M}}}
    \newcommand{\sO}{{\mathscr {O}}} \newcommand{\sP}{{\mathscr {P}}}
    \newcommand{\sQ}{{\mathscr {Q}}} \newcommand{\sR}{{\mathscr {R}}}
    \newcommand{\sS}{{\mathscr {S}}} \newcommand{\sT}{{\mathscr {T}}}
    \newcommand{\sU}{{\mathscr {U}}} \newcommand{\sV}{{\mathscr {V}}}
    \newcommand{\sW}{{\mathscr {W}}} \newcommand{\sX}{{\mathscr {X}}}
    \newcommand{\sY}{{\mathscr {Y}}} \newcommand{\sZ}{{\mathscr {Z}}}


% newcommand cal
    \newcommand{\CA}{{\mathcal {A}}} \newcommand{\CB}{{\mathcal {B}}}
    \newcommand{\CC}{{\mathcal {C}}} \newcommand{\CD}{{\mathcal {D}}}
    \newcommand{\CE}{{\mathcal {E}}} \newcommand{\CF}{{\mathcal {F}}}
    \newcommand{\CG}{{\mathcal {G}}} \newcommand{\CH}{{\mathcal {H}}}
    \newcommand{\CI}{{\mathcal {I}}} \newcommand{\CJ}{{\mathcal {J}}}
    \newcommand{\CK}{{\mathcal {K}}} \newcommand{\CL}{{\mathcal {L}}}
    \newcommand{\CM}{{\mathcal {M}}} \newcommand{\CN}{{\mathcal {N}}}
    \newcommand{\CO}{{\mathcal {O}}} \newcommand{\CP}{{\mathcal {P}}}
    \newcommand{\CQ}{{\mathcal {Q}}} \newcommand{\CR}{{\mathcal {R}}}
    \newcommand{\CS}{{\mathcal {S}}} \newcommand{\CT}{{\mathcal {T}}}
    \newcommand{\CU}{{\mathcal {U}}} \newcommand{\CV}{{\mathcal {V}}}
    \newcommand{\CW}{{\mathcal {W}}} \newcommand{\CX}{{\mathcal {X}}}
    \newcommand{\CY}{{\mathcal {Y}}} \newcommand{\CZ}{{\mathcal {Z}}}

    % newcommand frak
     \newcommand{\fa}{{\mathfrak{a}}}  \newcommand{\fb}{{\mathfrak{b}}}
     \newcommand{\fc}{{\mathfrak{c}}}  \newcommand{\fd}{{\mathfrak{d}}}
     \newcommand{\fe}{{\mathfrak{e}}}  \newcommand{\ff}{{\mathfrak{f}}}
     \newcommand{\fg}{{\mathfrak{g}}}  \newcommand{\fh}{{\mathfrak{h}}}
     \newcommand{\fii}{{\mathfrak{i}}}  \newcommand{\fj}{{\mathfrak{j}}}
     \newcommand{\fk}{{\mathfrak{m}}}  \newcommand{\fl}{{\mathfrak{l}}}
     \newcommand{\fm}{{\mathfrak{m}}}  \newcommand{\fn}{{\mathfrak{n}}}
     \newcommand{\fo}{{\mathfrak{o}}}  \newcommand{\fp}{{\mathfrak{p}}}
     \newcommand{\fq}{{\mathfrak{q}}}  \newcommand{\fr}{{\mathfrak{r}}}
     \newcommand{\fs}{{\mathfrak{s}}}  \newcommand{\ft}{{\mathfrak{t}}}
     \newcommand{\fu}{{\mathfrak{u}}}  \newcommand{\fv}{{\mathfrak{v}}}
     \newcommand{\fw}{{\mathfrak{w}}}  \newcommand{\fx}{{\mathfrak{x}}}
     \newcommand{\fy}{{\mathfrak{y}}}  \newcommand{\fz}{{\mathfrak{z}}}

    \newcommand{\fA}{{\mathfrak{A}}}  \newcommand{\fB}{{\mathfrak{B}}}
     \newcommand{\fC}{{\mathfrak{C}}}  \newcommand{\fD}{{\mathfrak{D}}}
     \newcommand{\fE}{{\mathfrak{E}}}  \newcommand{\fF}{{\mathfrak{F}}}
     \newcommand{\fG}{{\mathfrak{G}}}  \newcommand{\fH}{{\mathfrak{H}}}
     \newcommand{\fI}{{\mathfrak{I}}}  \newcommand{\fJ}{{\mathfrak{J}}}
     \newcommand{\fK}{{\mathfrak{K}}}  \newcommand{\fL}{{\mathfrak{L}}}
     \newcommand{\fM}{{\mathfrak{M}}}  \newcommand{\fN}{{\mathfrak{N}}}
     \newcommand{\fO}{{\mathfrak{O}}}  \newcommand{\fP}{{\mathfrak{P}}}
     \newcommand{\fQ}{{\mathfrak{Q}}}  \newcommand{\fR}{{\mathfrak{R}}}
     \newcommand{\fS}{{\mathfrak{S}}}  \newcommand{\fT}{{\mathfrak{T}}}
     \newcommand{\fU}{{\mathfrak{U}}}  \newcommand{\fV}{{\mathfrak{V}}}
     \newcommand{\fW}{{\mathfrak{W}}}  \newcommand{\fX}{{\mathfrak{X}}}
     \newcommand{\fY}{{\mathfrak{Y}}}  \newcommand{\fZ}{{\mathfrak{Z}}}



 % newcommand :rm
     \newcommand{\RA}{{\mathrm {A}}} \newcommand{\RB}{{\mathrm {B}}}
    \newcommand{\RC}{{\mathrm {C}}} \newcommand{\RD}{{\mathrm {D}}}
    \newcommand{\RE}{{\mathrm {E}}} \newcommand{\RF}{{\mathrm {F}}}
    \newcommand{\RG}{{\mathrm {G}}} \newcommand{\RH}{{\mathrm {H}}}
    \newcommand{\RI}{{\mathrm {I}}} \newcommand{\RJ}{{\mathrm {J}}}
    \newcommand{\RK}{{\mathrm {K}}} \newcommand{\RL}{{\mathrm {L}}}
    \newcommand{\RM}{{\mathrm {M}}} \newcommand{\RN}{{\mathrm {N}}}
    \newcommand{\RO}{{\mathrm {O}}} \newcommand{\RP}{{\mathrm {P}}}
    \newcommand{\RQ}{{\mathrm {Q}}} \newcommand{\RR}{{\mathrm {R}}}
    \newcommand{\RS}{{\mathrm {S}}} \newcommand{\RT}{{\mathrm {T}}}
    \newcommand{\RU}{{\mathrm {U}}} \newcommand{\RV}{{\mathrm {V}}}
    \newcommand{\RW}{{\mathrm {W}}} \newcommand{\RX}{{\mathrm {X}}}
    \newcommand{\RY}{{\mathrm {Y}}} \newcommand{\RZ}{{\mathrm {Z}}}

    \newcommand{\Ad}{{\mathrm{Ad}}} \newcommand{\Aut}{{\mathrm{Aut}}}
    \newcommand{\Br}{{\mathrm{Br}}} \newcommand{\Ch}{{\mathrm{Ch}}}
    \newcommand{\cod}{{\mathrm{cod}}} \newcommand{\cont}{{\mathrm{cont}}}
    \newcommand{\cl}{{\mathrm{cl}}}   \newcommand{\Cl}{{\mathrm{Cl}}}
    \newcommand{\disc}{{\mathrm{disc}}}\newcommand{\Eis}{{\mathrm{Eis}}}
    \newcommand{\Div}{{\mathrm{Div}}} \renewcommand{\div}{{\mathrm{div}}}
    \newcommand{\End}{{\mathrm{End}}} \newcommand{\Frob}{{\mathrm{Frob}}}
    \newcommand{\Gal}{{\mathrm{Gal}}} \newcommand{\GL}{{\mathrm{GL}}}
    \newcommand{\Hom}{{\mathrm{Hom}}} \renewcommand{\Im}{{\mathrm{Im}}}
    \newcommand{\Ind}{{\mathrm{Ind}}} \newcommand{\ind}{{\mathrm{ind}}}
    \newcommand{\inv}{{\mathrm{inv}}}
    \newcommand{\Isom}{{\mathrm{Isom}}} \newcommand{\Jac}{{\mathrm{Jac}}}
    \newcommand{\ad}{{\mathrm{ad}}}  \newcommand{\Tr}{{\mathrm{Tr}}}
    \newcommand{\Ker}{{\mathrm{Ker}}} \newcommand{\Ros}{{\mathrm{Ros}}}
    \newcommand{\Lie}{{\mathrm{Lie}}} \newcommand{\Hol}{{\mathrm{Hol}}}

    \newcommand{\cyc}{{\mathrm{cyc}}}\newcommand{\id}{{\mathrm{id}}}
    \newcommand{\new}{{\mathrm{new}}} \newcommand{\NS}{{\mathrm{NS}}}
    \newcommand{\ord}{{\mathrm{ord}}} \newcommand{\rank}{{\mathrm{rank}}}
    \newcommand{\PGL}{{\mathrm{PGL}}} \newcommand{\Pic}{\mathrm{Pic}}
    \newcommand{\cond}{\mathrm{cond}} \newcommand{\Is}{{\mathrm{Is}}}
    \renewcommand{\Re}{{\mathrm{Re}}} \newcommand{\reg}{{\mathrm{reg}}}
    \newcommand{\Res}{{\mathrm{Res}}} \newcommand{\Sel}{{\mathrm{Sel}}}
    \newcommand{\RTr}{{\mathrm{Tr}}} \newcommand{\alg}{{\mathrm{alg}}}
    \newcommand{\PSL}{{\mathrm{PSL}}}

\newcommand{\coker}{{\mathrm{coker}}}
\newcommand{\val}{{\mathrm{val}}} \newcommand{\sign}{{\mathrm{sign}}}
\newcommand{\mult}{{\mathrm{mult}}} \newcommand{\Vol}{{\mathrm{Vol}}}
\newcommand{\Meas}{{\mathrm{Meas}}}\renewcommand{\mod}{\ \mathrm{mod}\ }
\newcommand{\Ann}{\mathrm{Ann}}
\newcommand{\Tor}{\mathrm{Tor}}
\newcommand{\Supp}{\mathrm{Supp}}\newcommand{\supp}{\mathrm{supp}}
\newcommand{\Max}{\mathrm{Max}}
\newcommand{\Coker}{\mathrm{Coker}}
\newcommand{\Stab}{\mathrm{Stab}}
\newcommand{\Irr}{\mathrm{Irr}}\newcommand{\Inf}{\mathrm{Inf}}\newcommand{\Sup}{\mathrm{Sup}}
\newcommand{\rk}{\mathrm{rk}}\newcommand{\Fil}{\mathrm{Fil}}
\newcommand{\Sim}{{\mathrm{Sim}}} \newcommand{\SL}{{\mathrm{SL}}}
\newcommand{\Spec}{{\mathrm{Spec}}} \newcommand{\SO}{{\mathrm{SO}}}
\newcommand{\SU}{{\mathrm{SU}}} \newcommand{\Sym}{{\mathrm{Sym}}}
\newcommand{\sgn}{{\mathrm{sgn}}} \newcommand{\tr}{{\mathrm{tr}}}
\newcommand{\tor}{{\mathrm{tor}}}  \newcommand{\ur}{{\mathrm{ur}}}
\newcommand{\vol}{{\mathrm{vol}}}  \newcommand{\ab}{{\mathrm{ab}}}
\newcommand{\Sh}{{\mathrm{Sh}}} \newcommand{\Ell}{{\mathrm{Ell}}}
\newcommand{\Char}{{\mathrm{Char}}}\newcommand{\Tate}{{\mathrm{Tate}}}
\newcommand{\corank}{{\mathrm{corank}}} \newcommand{\Cond}{{\mathrm{Cond}}}
\newcommand{\Inn}{{\mathrm{Inn}}} \newcommand{\Spf}{{\mathrm{Spf}}}
\newcommand{\Mat}{{\mathrm{Mat}}}


    \font\cyr=wncyr10  \newcommand{\Sha}{\hbox{\cyr X}}
    \newcommand{\wt}{\widetilde} \newcommand{\wh}{\widehat} \newcommand{\ck}{\check}
    \newcommand{\pp}{\frac{\partial\bar\partial}{\pi i}}
    \newcommand{\pair}[1]{\langle {#1} \rangle}
    \newcommand{\wpair}[1]{\left\{{#1}\right\}}
    \newcommand{\intn}[1]{\left( {#1} \right)}
    \newcommand{\norm}[1]{\|{#1}\|}
    \newcommand{\sfrac}[2]{\left( \frac {#1}{#2}\right)}
    \newcommand{\ds}{\displaystyle}
    \newcommand{\ov}{\overline}
    \newcommand{\Gros}{Gr\"{o}ssencharaktere}
    \newcommand{\incl}{\hookrightarrow}
    \newcommand{\lra}{\longrightarrow}
     \newcommand{\ra}{\rightarrow}
    \newcommand{\imp}{\Longrightarrow}
    \newcommand{\lto}{\longmapsto}
    \newcommand{\bs}{\backslash}
    \newcommand{\nequiv}{\equiv\hspace{-7.8pt}/}
    \theoremstyle{plain}


\definecolor{energy}{RGB}{114,0,172}
\definecolor{freq}{RGB}{45,177,93}
\definecolor{spin}{RGB}{251,0,29}
\definecolor{signal}{RGB}{203,23,206}
\definecolor{circle}{RGB}{217,86,16}
\definecolor{average}{RGB}{203,23,206}
\newcommand{\K}{\operatornamewithlimits{K}}
\colorlet{shadecolor}{gray!20}
\pgfplotsset{compat=1.9}
\def\N{10}
\def\M{4}
\usepgflibrary{fpu}

\algnewcommand\algorithmicinput{\textbf{Input:}}
\algnewcommand\algorithmicoutput{\textbf{Output:}}

\algnewcommand\Input{\item[\algorithmicinput]}
\algnewcommand\Output{\item[\algorithmicoutput]}

\def\upint{\mathchoice%
    {\mkern13mu\overline{\vphantom{\intop}\mkern7mu}\mkern-20mu}%
    {\mkern7mu\overline{\vphantom{\intop}\mkern7mu}\mkern-14mu}%
    {\mkern7mu\overline{\vphantom{\intop}\mkern7mu}\mkern-14mu}%
    {\mkern7mu\overline{\vphantom{\intop}\mkern7mu}\mkern-14mu}%
  \int}
\def\lowint{\mkern3mu\underline{\vphantom{\intop}\mkern7mu}\mkern-10mu\int}



\makeatletter
\let\c@equation\c@thm
\raggedbottom
\makeatother
\numberwithin{equation}{section}
%--------Meta Data: Fill in your info------
\author{Hojune Lee, 20210541}

\title{Abstract Algebra}
\begin{document}

\begin{titlepage}
    \begin{center}
        \vspace*{9cm}
            
        \Huge
        \textbf{Algorithms for NP-Hard Problems}
    
        \vspace{1cm}
        \large
        Summary note, Lectured by Eunjung Kim, KAIST
        \vspace{3cm}
        
        \LARGE
        \textbf{Hojune Lee}
            
        \vspace{8cm}
            
        \normalsize
        \textbf{School of Computing, KAIST}\\  
    \end{center}
\end{titlepage}

\hypersetup{linkcolor=black}
\tableofcontents
\hypersetup{linkcolor=blue}

\newpage 

\chapter{Parameterized Algorithm}

\begin{tcolorbox}[colback=yellow!10!white,colframe=brown!75!black,title=Parameterization]

The \textit{parameterization} is a map defined by 
\[\kappa : \Sigma^* \rightarrow \BN_0\]
With this definition, \textit{parameterized problem} is an ordered pair decision problem for language $\fL \subseteq \Sigma^*$ and parameterization $\kappa$, 
$(\fL, \kappa)$. Since $\kappa$ is a map, we can take any well-defined function for $\kappa$. So, we can claim that for wrong encoded 
instance in $\Sigma^*$, they goes to some fixed value, e.g. $0$. 

Here, it is important that why we take this definition. I'll write the intuition later, after introducing the example. 
\end{tcolorbox}

\begin{tcolorbox}[colback=yellow!10!white,colframe=red!75!black,title=\textsc{Vertex Cover}]

All problems are defined in red-colored box. First problem is \textsc{Vertex Cover}. For a graph $G = (V, E)$ and $k \in \BN_0$, 
\textit{vertex cover} is a set $S \subseteq V$ such that $\forall uv \in E, u \in S \vee v \in S$. i.e. All edges in $E$ are covered by $S$. 

Here, \textsc{Vertex Cover} is a decision problem of following encoded language $\fL$ : 
\[\fL_{VC} = \{(G, k) | G = (V, E), k \in \BN_0, G \text{ has vertex cover whose size at most } k\}\]

\end{tcolorbox}

Since $(G, k) \in \Sigma^*$ is a encoded string, they are together. Means that, when we describe the input size $|x|$ of algorithm, 
it doesn't contain the seperated meaning of size of $G$ and magnitude of $k$. So we parameterization $\kappa$ here, to extract the meaning of $k$. 
\[\kappa : (G, k) \mapsto k\] 
As we'll see later, another information also can be extracted : 
\[\kappa : (G, k) \mapsto k - lp^*(G)\]
This is conceptual meaning of parameterization of a problem. 

\begin{tcolorbox}[colback=yellow!10!white,colframe=brown!75!black,title=Fixed Parameter Tractable]

A parameterized problem $(\fL, \kappa)$ is \textit{fixed parameter tractable} if there is an algorithm such that 
for every $x \in \Sigma^*$, decides whether $x \in \fL$ or not in time $f(\kappa(|x|))\cdot |x|^{\CO(1)}$

\end{tcolorbox}

\newpage 

\chapter{Branching Algorithms}

\section{$\CO(2^k)$-time Algorithm for Vertex Cover}

\begin{algorithm}

\caption{$\CO(2^k)$-time Algorithm for Vertex Cover}

\begin{algorithmic}[1]
\Procedure{\textbf{VC}}{$G,k$}
\Input{$G = (V, E)$ is a graph, $k \in \BN_0$}
\Output{Decide $(G, k) \in \fL_{VC}$ or not }
\If{$|E| = 0$}
    \Return Yes
\ElsIf{$k = 0$}
    \Return No 
\Else   
\State $uv \gets Pick(E)$ \Comment{Pick any edge}
\State \Return \textbf{VC}$(G - u, k - 1)$ $\vee$ \textbf{VC}$(G - v, k - 1)$
\EndIf
\EndProcedure
\end{algorithmic}

\end{algorithm}

This is very fundamental algorithm for \textsc{Vertex Cover}. We can prove following: 

\begin{tcolorbox}[colback=yellow!10!white,colframe=green!75!black,title=Algorithm 1]

Algorithm 1 is \textit{correct}(i.e., decides whether $(G, k) \in \fL_{VC}$ or not correctly) in time $\CO(2^k \cdot poly(n))$

\end{tcolorbox}

\begin{proof}
    
We can say that algorithm 1 is correct means that 
\[\mathbf{VC}(G, k) = \text{Yes} \iff (G, k) \in \fL_{VC}\]
We'll prove via mathematical induction on $k \in \BN_0$. If $k = 0$, it is trivial.
Then assume that for any $G = (V, E)$, $VC(G, i) = \text{Yes} \iff (G, i) \in \fL_{VC}, 0 \leq i < k$. 
And one easy fact is that, if $S$ is a vertex cover of $G$, and $u \in S$, then $S \setminus \{u\}$ is a vertex cover of $G - u$. 
($\Leftarrow$) Let $(G, k) \in fL_{VC}$. Then, exists a vertex cover $S$ of $G$ such that $|S| \leq k$. 
Suppose that given instance is non-trivial, so reach at line 4. Then pick $\forall uv \in E$, and W.L.O.G., $u \in S$. Then by fact, 
$S \setminus \{u\}$ is a vertex cover of $G - u$ whose size is at most $k - 1$. Then, we can determine that $\mathbf{VC}(G-u,k-1) = \text{Yes}$ by hypothesis. 
\newline
($\Rightarrow$) Suppose that $\textbf{VC}(G, k) = \text{Yes}$. Then for non-trivial cases, W.L.O.G.,
$\textbf{VC}(G-u, k-1)$ is Yes. It implies that $(G-u, k-1) \in \fL_{VC}$. Then, let $S'$ be the vertex cover of $G-u$ whose size is at most $k-1$. 
Then, let $S = S' \cup \{u\}$. It is clearly a vertex cover of $G$ whose size is at most $k$. So $(G, k) \in \fL_{VC}$. 

Now, for the running time. It is obvious that line 2,3,5 effects polynomial time w.r.t. input size. And there are at most $k$-depth of binary recursion, so 
running time is $\CO(2^k \cdot poly(n))$. 

\end{proof}

\section{$\CO((3k)^k)$-time Algorithm for Feedback Vertex Set}

Now we define next problem to solve, which is called \textsc{Feedback Vertex Set}. 

\begin{tcolorbox}[colback=yellow!10!white,colframe=red!75!black,title=\textsc{Feedback Vertex Set}]

For any graph $G = (V, E)$, a \textit{feed-back vertex set} of $G$ is a set $S \subseteq V$ such that 
\[\forall C \in Cycle(G), S \cap C \neq \emptyset\]
i.e., set of vertices that intersect with all of cycles in $G$. So, we can define decision problem \textsc{Feedback Vertex Set} by 
\[\fL_{FVS} = \{(G, k) | G = (V, E), k \in \BN_0, G \text{ has FVS whose size at most $k$}\}\]

\end{tcolorbox}

\begin{algorithm}
    
\caption{$\CO((3k)^k)$-time Algorithm for Feedback Vertex Set}

\begin{algorithmic}[1]
\Procedure{\textbf{FVS}}{$G,k$}
\If{$|E| = 0$}
    \Return Yes
\EndIf
\If{$k = 0$}
    \Return No
\EndIf
\While{$\exists v \in V(G), d_G(v) \leq 2$}
\If{$d_G(v) \leq 1$}
    \State $G \gets G - v$
\ElsIf{$|N_G(v)\setminus \{v\}| = 2$}
    \State $\{x, y\} \gets N_G(v) \setminus \{v\}$
    \State $G \gets G - v + xy$ \Comment{Allow multigraphs}
\Else    
    \State $G \gets G \setminus N_G(v)$
    \State $k \gets k - 1$ \Comment{Assume that we choose the adjacent vertex of $v$}
\EndIf 
\EndWhile
\State $\{v_n, \cdots, v_1\} \gets Sort(V(G))$ \Comment{Note: Decreasing Order}
\State \Return $\bigvee_{i=1}^{3k} \mathbf{FVS}(G - v_i, k-1)$
\EndProcedure
\end{algorithmic}

\end{algorithm}

\begin{tcolorbox}[colback=yellow!10!white,colframe=green!75!black,title=Algorithm 2]

Algorithm 2 is correct for \textsc{Feedback Vertex Set} in time $\CO((3k)^k \cdot poly(n))$

\end{tcolorbox}

\begin{lem}
    
For any feed-back vertex set $X$ of $G = (V, E)$ holds that 
\[\sum_{x \in X} (d_G(x) - 1) \geq |E| - |V| + 1\]

\end{lem}

\begin{proof}
    
One easily obtained fact is that, $G[V \setminus X]$ is a forest. Means that, 
\[|E(G[V\setminus X])| \leq |V(G[V\setminus X])| - 1 = |V| - |X| - 1\]
However, 
\[|E| = \underbrace{|E(G[X])| + |E(X, V \setminus X)|}_{\leq \sum_{x \in X} d_G(x)} + |E(G[V \setminus X])|\]
So, 
\[\sum_{x \in X}d_G(x) \geq |E| - |E(G[V \setminus X])| \geq |E| - |V| + |X| + 1\]
Implies that 
\[\sum_{x \in X} (d_G(x) - 1) \geq |E| - |V| + 1\]

\end{proof}

\begin{lem}
    
Let $G = (V, E)$ such that $\delta(G) \geq 3$. Let $V = \{v_1, v_2, \cdots, v_n\}$ where 
\[d_G(v_1) \geq d_G(v_2) \geq \cdots \geq d_G(v_n)\]
Then, for any feedback vertex set $X$ of $G$, 
\[|X| \leq k \implies X \cap \{v_1, \cdots, v_{3k}\} \neq \emptyset\]

\end{lem}

\begin{proof}
    
Suppose not. By hand-shake lemma, 
\[\sum_{v \in V} d_G(v) = 2|E| \geq 3|V| \qquad \text{$\because \delta(G) \geq 3$}\]
And since $X \cap \{v_1, \cdots, v_{3k}\} = \emptyset$, 
We can obtain following from above lemma:
\[\sum_{i = 1}^{3k} (d_G(v_i)-1) \geq 3 \sum_{x \in X} (d_G(x)-1) \geq 3|E| - 3|V| + 3 \] 
\[\sum_{i = 3k+1}^{n} (d_G(v_i) - 1) \geq \sum_{x \in X} (d_G(x) - 1) \geq |E| - |V| + 1\]
implies that, 
\[\sum_{v \in V(G)} (d_G(v) - 1) = 2|E| - |V| \geq 4|E| - 4|V| + 4\]
It deduce 
\[3|V| > 2|E|\]
which contradict to first line.
\end{proof}

With second lemma, one can show that Algorithm 2 is correct. And since it generates at most $3k$ branches for each step, and the maximal depth is $k$, so the running time will be $\CO((3k)^k \cdot poly(n))$

\section{$\CO(4^{k-lp^*(G)})$-time Algorithm for Vertex Cover}

In this section, we review the \textsc{Vertex Cover} with another parameterization, which is 
\[\kappa : (G, k) \mapsto k - lp^*(G)\]
Here, $lp^*(G)$ is a solution of another problem, which called linear program version of vertex cover. 
\begin{tcolorbox}[colback=yellow!10!white,colframe=red!75!black,title=Linear Program for Vertex Cover]

For $G = (V, E)$, linear program for vertex cover is called \textsc{LPVC}, which defined as following. 
\[\min \sum_{u \in V(G)} x_u\]
where 
\[\forall uv \in E, x_u + x_v \geq 1, \qquad  \forall u \in V, x_u \geq 0 \tag{$\star$}\]
holds. Let's call above condition by $\star$. And we denote $lp^*(G)$ be the optimal solution of \textsc{LPVC}. 

\end{tcolorbox}

One fact is that, 
\[\text{size of optimal vertex cover of $G$} \geq lp^*(G)\]
This is because that, optimal vertex cover is very special case of \textsc{LPVC}. It can be easily obtain, since 
every vertex cover satisfy the condition $\star$. So, this automatically implies that 
for given $(G, k)$, if $k < lp^*(G)$, $(G, k)$ is No instance of \textsc{Vertex Cover}. 

Now, we'll explore the algorithm which related to these LP problems. Before move on, we'll take a following theorem without proof. 

\begin{tcolorbox}[colback=yellow!10!white,colframe=blue!75!black,title=Linear Program is P]

Linear Program can be solved in polynomial time. So that, \textsc{LPVC} also solved in polynomial time.

\end{tcolorbox}

And now we'll define half-integral solution for \textsc{LPVC}. 

\begin{tcolorbox}[colback=yellow!10!white,colframe=brown!75!black,title=Half Integral solution for LPVC]

We call that the solution $x^*$ for \textsc{LPVC} is \textit{half-integral} when 
\[\forall u \in V(G), x^*_u \in \{0, \frac{1}{2}, 1\}\]

\end{tcolorbox}

With the above fact, we can conduct following decomposition in polynomial time. 

\begin{tcolorbox}[colback=yellow!10!white,colframe=green!75!black,title=Crown Decomposition]

Let $x^*$ be the solution of \textsc{LPVC}. Then we can partition $V(G)$ into following 3 sets : 
\[C = \{u \in V(G) | x^*_u < 0.5\}\]
\[H = \{u \in V(G) | x^*_u > 0.5\}\]
\[R = \{u \in V(G) | x^*_u = 0.5\}\]
This is called \textit{crown decomposition} of $G$. This decomposition can be obtained in polynomial time. 
\end{tcolorbox}
\begin{figure}[htbp]
    \centering
    \includegraphics[width=0.7\textwidth]{crown-decomposition.png}
    \caption{A graph illustrating the crown decomposition.}
    \label{fig:crown-decomposition}
\end{figure}

Now, we are going to argue many properties on crown decomposition. 

\begin{lem}

For every $H' \subseteq H$, $|N_G(H') \cap C| \geq |H'|$

\end{lem}

\begin{proof}
    
Actually, this is quite intuitive. If not, that situation is very un-natural to reduce the total-sum. 

Suppose not, i.e., there exists $H' \subseteq H$ such that $|N_G(H') \cap C| < |H'|$. 
Now idea is very simple, let's construct more cheap solution from here! 
We can define following value $\epsilon > 0$ 
\[\epsilon = \min_{h \in H'} \{x^*_h - 0.5\} > 0\]
And then, we'll construct more cheap solution by 
\[\forall u \in V(G), x'_u = \begin{cases}
    x^*_u + \epsilon & u \in N_G(H') \cap C \\
    x^*_u - \epsilon & u \in H' \\
    x^*_u & \text{otherwise}
\end{cases}\]

$x'$ is feasible LP solution since all our reponse to check is for incident parts of $H'$, but it preserves 
$H'$ to $C$ sum, and since $x'$ have same crown decomposition, $H'$ to $H$ and $H'$ to $R$ also resolved. However, as 
$|N_G(H') \cap C| < |H'|$, the total cost $|x'|$ is strictly small than $|x^*|$. 
\[\sum_{u \in V(G)} x'_u = \left(\sum_{u \in V(G)} x^*_u \right) - \epsilon(|H'| - |N_G(H') \cap C|) < \sum_{u \in V(G)} x^*_u\]
This contradicts to $x^*$ is an optimal solution. 

\end{proof}

\begin{lem}
    
Let $(R, H, C)$ be a crown decomposition of $G$ w.r.t. optimal solution $x^*$ of \textsc{LPVC}. 
Then, following solution is also a optimal solution of \textsc{LPVC}. 

\[x'_u = \begin{cases}
    0.5 & u \in R \\
    1 & u \in H \\
    0 & u \in C
\end{cases}\]

\end{lem}

\begin{proof}
    
Let $\CX$ be the set of all optimal solution of \textsc{LPVC} such that yields crown decomposition $(R, H, C)$. 
And for each $x^* \in \CX$, we can define following two sets: 
\[C_{x^*} = \{c | c \in C , x^*_c > 0\}\]
\[H_{x^*} = \{h | h \in H, x^*_h < 1\}\]
Then we can change our claim into that there exists $x' \in \CX$ such that $C_{x'} \cap H_{x'} = \emptyset$. 
Suppose not, and let $x'$ be the solution that minimize $|C_{x'} \cap H_{x'}| > 0$. In this case, we can choose following $\delta > 0$. 
\[\delta = \min \left( \min_{c \in C_{x'}} \{x'_c\}, \min_{h \in H_{x'}} \{1 - x'_h\} \right)\]
Then here, we can construct following revised solution $x''$ by 
\[x''_u = \begin{cases}
    x'_u - \delta & u \in H_{x'} \\
    x'_u + \delta & u \in C_{x'} \\
    x'_u & \text{otherwise}
\end{cases}\]
It is obvious that $x''$ is feasible solution of \textsc{LPVC}. Actually, $x''$ is an optimal solution. 
First, for all vertices in $H_{x'}$ can't be incident with $C \setminus C_{x'}$. (Because, vertices in $C_{x'}$ assigned $0$)
It implies that, $N_G(H_{x'}) \cap C \subseteq C_{x'}$. However, we know that $|H_{x'}| \leq |N_G(H_{x'}) \cap C|$ by above lemma. 
So finally, we these two relation, 
\[|H_{x'}| \leq |C_{x'}| \implies |x''| \leq |x'|\]
Since $x'$ is optimal, $|x''| = |x'|$ and so $x''$ also optimal solution. However, at least one vertex in $H_{x'}$ or $C_{x'}$ are collapsed into $1$, $0$ respectively, 
so $|C_{x''} \cap H_{x''}| < |C_{x'} \cap H_{x'}$ which contradicts to choice of $x'$. So that, there exists $x' \in \CX$ such that 
$C_{x'} \cap H_{x'} = \emptyset$, holds our claim. 
\end{proof}

\textbf{Remark} During the proof, we construct the algorithmic way to convert any optimal LPVC solutions into half-integral optimal LPVC solution. 
Since we can obtain any LPVC solution in polynomial time and above algorithm also, 
we can find half-integral optimal LPVC solution $x^*$ in polynomial time always. 

\begin{tcolorbox}[colback=yellow!10!white,colframe=green!75!black,title=Obtain Half-Integral Solution]

For given graph $G = (V, E)$, we can obtain half-integral optimal LPVC solution $x^*$ in polynomial time.
And define 
\[V_a(x^*) = \{v \in V | x^*_v = a\}\]
Then here, crown decomposition w.r.t. $x^*$ is $(V_{0.5}, V_1, V_0)$. 
\end{tcolorbox}

\begin{lem}
    
For $G' = G[V_{0.5}]$, $\forall v \in G', x'_v = 0.5$ is an optimal solution of LPVC of $G'$.

\end{lem}

\begin{proof}
    
Very easy. We can observe that, on the crown decomposition of $G$, all $V_{0.5}$ vertices are incident with $V_{0.5}$ itself or $V_1$.
So, if there exists another optimal solution than $x'$, we can transplant that solution directly in $G$, then we can make more optimal solution to LPVC of $G$. 
This occur contradiction, so $x'$ is an optimal solution of LPVC of $G[V_{0.5}]$. 

\end{proof}

\begin{lem}

There exists an optimal vertex cover containing $V_1$. 

\end{lem}

\begin{proof}
    
Let $S$ be any optimal vertex cover of $G$. Then, w.r.t. the crown decomposition of $G$ by $(V_{0.5}, V_1, V_0)$, 
we can define 
\[S_0 = S \cap V_0, \quad S_{0.5} = S \cap V_{0.5}, \quad S_1 = S \cap V_1\]
Then let's construct new set 
\[S' = (S \setminus S_0) \cup V_1\]
i.e., replace all cover vertices in $C$ into remaining vertices in $H$. One can be easily observed, $S'$ is a vertex cover of $G$. 
So, to check $S'$ also an optimal vertex cover, we just need to check that number of removed vertices ($|S_0|$) is larger or equal than number of added vertices ($|V_1 \setminus S_1|$). 
Not too hard to see. Imagine the crown decomposition diagram. Since $S$ was vertex cover, all the edges are covered. It means that, there aren't any edges between 
$V_0 \setminus S_0$ and $V_1 \setminus S_1$. So, $N(V_1 \setminus S_1) \cap V_0 \subseteq S_0$. 
And according to the Lemma 2.3.1., 
\[|V_1 \setminus S_1| \leq |N(V_1 \setminus S_1) \cap V_0| \leq |S_0|\]
This completes the proof.

\end{proof}

\begin{lem}
    
The instance $I' = (G', k')$ is equivalent to $I = (G, k)$ where $G' = G[V_{0.5}]$ and $k' = k - |V_1|$. 
Furthermore, it holds that $k' - lp^*(G[V_{0.5}]) = k - lp^*(G)$

\end{lem}

\begin{proof}
    
Let's prove that $I' \iff I$ first. 

($\Rightarrow$) 
Suppose that $(G', k')$ is Yes instance. Then, take a vertex cover $S'$ for $G'$ whose size is at most $k'$. Then, for $G$, just set $S = S' \cup V_1$. 
It is clearly vertex cover of $G$ and \[|S| = |S'| + |V_1| \leq k' + |V_1| = k\] so $(G, k)$ is Yes instance. 

($\Leftarrow$)
Suppose that $(G, k)$ is Yes instance. Then, by Leamma 2.3.4. there exists an optimal vertex cover $S$ containing $V_1$. 
Then, $S_{1/2} = S \setminus V_1 \subseteq V_{1/2}$ holds. Then, this cover all the vertices in $V_{1/2} \setminus S_{1/2}$ since 
if not, it is contradicts to optimal LPVC solution. So, $S_{1/2}$ is vertex cover of $G[V_{1/2}]$. However, since $|S| \leq k$, \[|S_{1/2}| = |S| - |V_1| \leq k - |V_1| = k'\]
So, $(G', k')$ is Yes instance. 

For the further statement, we already have lemma that $lp^*(G[V_{0.5}]) = \frac{1}{2}|V_{1/2}|$. 
And for optimal LPVC half-integral solution, $lp^*(G) = |V_1| + \frac{1}{2}|V_{1/2}|$. So the claim holds. 
\end{proof}

\textbf{Remark} We show that $\forall v \in V_{1/2}, x^*_v = 1/2$ can be optimal LPVC solution for $G[V_{1/2}]$. It does not guarantee that 
it is unique. 

\vspace{4mm}

Now, with Lemma 2.3.5. we can reduce(or preprocess) our original instance $(G, k)$. 

\begin{tcolorbox}[colback=yellow!10!white,colframe=blue!75!black,title=Reduction Rule of LPVC]

Let $(G, k)$ be given instance. 
\begin{enumerate}
    \item For every half-integral optimal LPVC solution $y^*$ of $(G, k)$,  
    \item If there aren't any vertex $v$ such that $y^*_v = 1$, then end. 
    \item If such $v$ exists, construct $(G[V_{1/2}(y^*)], k - |V_1(y^*)|)$
\end{enumerate}
It holds because of Lemma 2.3.5.
\end{tcolorbox}

With apply forever this rule, we finally obtain $(G', k')$ whose optimal half-integral LPVC solution is 
$\forall v \in V(G'), x^*_v = 1/2$. Then, we can construct new \textsc{Vertex Cover} algorithm. 

\begin{algorithm}
    
\caption{$\CO(4^{k-lp^*(G)})$-time algorithm for Vertex Cover}

\begin{algorithmic}[1]
\Procedure{\textbf{VC-above-LP}}{$G, k$}
\State $(G, k) \gets$ Apply reduction rule until converge
\State // Here, all-$\frac{1}{2}$ is the unique half-integral optimal solution. 
\If{$k < lp^*(G) = \frac{1}{2}|V(G)|$} 
    \State \Return No \Comment{We observed this fact}
\ElsIf{$k \geq 2 \cdot lp^*(G) = |V(G)|$}
    \State \Return Yes \Comment{Because we can pick every vertices}
\Else   
    \State $uv \gets Pick(E(G))$
    \State \Return \textbf{VC-above-LP}$(G-u, k-1)$ $\vee$ \textbf{VC-above-LP}$(G-v, k-1)$
\EndIf
\EndProcedure
\end{algorithmic}

\end{algorithm}

It is easy to see that, this algorithm works correctly. Let's consider following measure: 
\[\mu(G, k) = k - lp^*(G)\]
i.e., this is our new parameterization for Vertex Cover above LP. Then, our goal is that, 
by showing this factor decreases at least $\frac{1}{2}$ per each step, so our algorithm terminated and 
depth of tree is at most $2(k - lp^*(G))$. 

\begin{prop}
For every iteration, the measure $\mu$ decreases at least $\frac{1}{2}$. 
\end{prop}

\begin{proof}
    
Here we claim that, for every $v \in V(G)$, 
\[lp^*(G - v) \geq lp^*(G) - \frac{1}{2}\]
Suppose not. Then, $lp^*(G - v) < lp^*(G) - \frac{1}{2}$. Since each $lp^*$ value is quantumized by $1/2$, it implies that 
\[lp^*(G-v) \leq lp^*(G) - 1\]
Then, let's choose the half-integral optimal LPVC solution $x^*$ of $G - v$. 
Then, append $x^*$ for $v$ with $x^*_v = 1$. Then such new $x^*$ holds size $lp^*(G)$ optimal solution of LPVC of $G$, however, 
all-$1/2$ is unique solution. It is contradiction, so claim holds. 
With that, for every iteration, W.L.O.G. for vertex $v$,
\[\mu(G-v, k-1) = k - 1 - lp^*(G-v) \leq k - lp^*(G) - \frac{1}{2} = \mu(G, k) - \frac{1}{2}\]
However, our algorithm guarantees that if $k < lp^*(G)$ or $k \geq 2 \cdot lp^*(G)$ then terminate. 
So, if $k - lp^*(G) < 0$, then terminate. So we can decrease our measure at most $2\mu(G, k)$ times. 
So the running time is $\CO(4^{k - lp^*(G)})$. And also, we guaranteed that $k - lp^*(G) \leq k/2$. So, it works as well as $\CO(2^k)$ time. 

\end{proof}

\newpage

\chapter{Kernelization}

\section{$\CO(k^2)$-Vertex Kernel}

\begin{tcolorbox}[colback=yellow!10!white,colframe=brown!75!black,title=Kernelization]

A \textit{kernelization} of problem $P$ is a polynomial time algorithm (w.r.t. $|x|$ and $k$ both) 
such that transform any instance $(x, k)$ of $P$ into another instance $(x', k')$ of $P$ satisfying 
\begin{itemize}
    \item equivalence: $(x, k) \in P$ if and only if $(x', k') \in P$
    \item size bound: $|x'| \leq g(k)$ and $k' \leq g(k)$ for some computable function $g$. 
\end{itemize}
The final instance $(x', k')$ is called \textit{kernel}, and $g(k)$ be the size of a kernel. 

\end{tcolorbox}

\begin{tcolorbox}[colback=yellow!10!white,colframe=brown!75!black,title=Reduction Rules]

Kernelization consists in applying some package of rules. A \textit{reduction rules} for $P$ is a 
polynomial time algorithm (w.r.t. $|x|$ and $k$ both) such that transfrom any instance $(x, k)$ of $P$ into another instance $(x', k')$ of $P$ satisfying 
equivalence. (Equivalence is also called by \textit{safeness, soundness}) We say that a instance $(x, k)$ is \textit{irreducible} if 
reduction rule $R$ does not change them. 

\end{tcolorbox}

Intuitively, if we can reduce the instance size as much as we want (bounded by $g(k)$) in polynomial time, 
then we can imagine that solving time of the reduced problem only depends on $g(k)$, the new input size. Following lemma is based on this intuition. 

\begin{lem}
    
Let $P$ be the parameterized decidable problem. $P$ is fixed parameter tractable (FPT) if and only if it admits a kernelization. 

\end{lem}

\begin{proof}
    
($\Leftarrow$) Suppose that we can find kernelization: 
\[(x, k) \mapsto (x', k') \quad |x'|, k' \leq g(k)\]
Then since $|x'| \leq g(k)$, we can solve this problem by brute-force in $h(g(k))$ time. So this is FPT. 

($\Rightarrow$) Suppose that $P$ is FPT. i.e., we can decide $P$ in time $f(k) \cdot |x|^c$ time. Then, running this algorithm for time $|x|^{c+1}$. 
If it terminates, then return any small-sized constant instance. Otherwise, it means that $|x| < f(k)$. It directly means that given original instance is a kernel already. 

\end{proof}

Now, let's see the first example. 

\begin{tcolorbox}[colback=yellow!10!white,colframe=red!75!black,title=Buzz' Kernelization: $\CO(k^2)$-Vertex Kernel]

This is kernelization algorithm for \textsc{Vertex Cover}. It input $(G, k)$ and output $(G', k')$ such that 
$|V(G')| \leq \CO(k^2)$ and $k' \leq k$ with safeness. (So it is valid kernelization). 

\begin{itemize}
    \item Reduction Rule 1: $\exists v \in V(G), d_G(v) = 0 \implies (G, k) \mapsto (G - v, k)$
    \item Reduction Rule 2: $\exists v \in V(G), d_G(v) \geq k + 1 \implies (G, k) \mapsto (G-v, k-1)$
\end{itemize}

\end{tcolorbox}

It is easy to see that Rule 1 and 2 are safe. (i.e., after any reduction, $(G, k)$ is yes instance if and only if $(G', k')$ is yes instance.) 

Now, let's analyze the size of kernel. Let $(G', k')$ be the kernel(maybe) of above reduction rules. 
Then, $k' \leq k$ is clear. The measurement of size of kernel is quite tricky. We claim that if $(G, k)$ is irreducible w.r.t. above rules, then $|G|(= |V(G)) \leq \CO(k^2)$. 
\begin{proof}

Suppose that $(G, k)$ is yes instance. Then there exists a vertex cover $C$ of $G$. Here, 
\[|C| \leq k ,\quad \forall v \in V(G), 1 \leq d_G(v) \leq k \implies |E(C)| \leq k^2\]
However, $V(G) \setminus C$ is independent set in $G$. And there aren't any isolated vertices in $G$, 
$|V(G) \setminus C| \leq |E(C)| \leq k^2$. So, 
\[|G| = |V(G)| = |C| + |V(G) \setminus C| \leq k + k^2 \in \CO(k^2)\]

Next one is tricky. If $|G| > \CO(k^2)$, it immediately implies that $(G, k)$ is no-instance. So, in this case, we can 
output tiny no-instance example for kernel output. 

\end{proof}

\newpage 

\section{LP-based kernelization for Vertex Cover}
Actually, this is sort of toy example. Now we'll see amazing results further. 
Before move on, let us enable to use some useful decomposition of graphs without proof. Actually, we already see this in previous chapter. 

\begin{tcolorbox}[colback=yellow!10!white,colframe=red!75!black,title=Nemhauser-Trotter Theorem,breakable]

Given any graph $G$, we can find partition $V(G) = (R, H, C)$ holds followings in polynomial time. 
\begin{itemize}
    \item For any vertex cover $S_r$ of $G[R]$, $S_r \cup H$ is a vertex cover of $G$. 
    \item There exists an optimal vertex cover containing $H$. 
    \item Any vertex cover of $G[R]$ is of size at least $\frac{1}{2}|R|$. 
\end{itemize}

\end{tcolorbox}

Then, let's define kernelization related to this theorem. 

\begin{tcolorbox}[colback=yellow!10!white,colframe=red!75!black,title=$\CO(2k)$-kernelization of Vertex Cover]

This is kernelization algorithm for \textsc{Vertex Cover}. It takes input $(G, k)$ and output $(G', k')$ such that 
$|V(G)| \leq 2k$. 

\begin{itemize}
    \item Rule 1: $V(G) \mapsto (R, H, C)$, $H \cup C \neq \emptyset \implies (G, k) \mapsto (G[R], k - |H|)$
\end{itemize}

\end{tcolorbox}

\begin{lem}
    This holds safeness. 
\end{lem}

\begin{proof}
    
($\Leftarrow$) Let $S'$ be a vertex cover of $G[R]$ whose size is at most $k - |H|$. Then, by the condition 1 of NT theorem, 
$S = S' \cup H$ be the vertex cover of $G$ whose size is at most $k$. 

($\Rightarrow$) Let $S$ be a vertex cover of $G$ whose size is at most $k$. Then, there exists a vertex cover $S$ of $G$, such that $H \subseteq S$. 
One can be easily derived from condition 1: $C$ is independent set of $G$ and $R$ does not incident with $C$. First one is easy, and for second one, suppose that $cr \in E(G)$ for $c \in C, r \in R$. Then, we can choose $R - r$ for the vertex cover of $G[R]$ obviously. However, 
since $H \cup (R - r)$ does not cover $cr$, it is contradiction. So, $R$ does not incident with $C$. So that, $H$ cover all the incident edges of $C$. 
Since $S$ is optimal, $S \setminus H$ is vertex cover of $G[R]$ whose size is at most $k - |H|$. 

\end{proof}

\begin{lem}
    
The kernel size($=|V(G)|$) of above kernelization is $2k$. 

\end{lem}

\begin{proof}
    
Let $(G, k)$ be irreducible instance of above algorithm. Then, it implies that for NT decomposition $(R, H, C)$ of $G$, 
$H \cup C = \emptyset$. Then, it means that $G[R] = G$. So, by condition 3, any vertex cover of $G[R] = G$ is size at least $\frac{1}{2}|R| = \frac{1}{2}|V|$. 
So, if $k < \frac{1}{2}|V|$, it automatically becomes no instance. So we can output tiny no-instance kernel for this case, $|V| > 2k$. Otherwise, 
$|V| \leq 2k$. So kernel size of $(G, k)$ is less than $2k$. 

\end{proof}

\newpage 

\section{$3k$-vertex kernel using Crown Decomposition}

\begin{tcolorbox}[colback=yellow!10!white,colframe=brown!75!black,title=Crown Decomposition]

A \textit{crown decomposition} of a graph $G = (V, E)$ is a partition $(R, H, C)$ of $V$ s.t. 
\begin{itemize}
    \item $H \cup C \neq \emptyset$
    \item $C$ is independent set of $G$
    \item $N(C) \subseteq H$, i.e., $H$ seperates $C$ and $R$. 
    \item $\forall H' \subseteq H$, $|N(H')\cap C| \geq |H'|$, i.e., Hall's condition holds between $H, C$
\end{itemize}

\end{tcolorbox}

\begin{tcolorbox}[colback=yellow!10!white,colframe=brown!75!black,title=Matching and Saturated]

A \textit{matching} in a graph $G = (V, E)$ is a set of disjoint edges. We call that 
$M$ is maximal matching when if we add any edge to $M$, it is not matching anymore. We call that $M$ is maximum matching when 
$M$ is the largest matching. We say that $v \in V(G)$ is \textit{saturated} by matching $M$ if $v$ is incident with $M$. 

\end{tcolorbox}

We have two famous theorems in Graph Theory course, which are Hall's theorem and k\"onig's theorem. 

\begin{tcolorbox}[colback=yellow!10!white,colframe=red!75!black,title=Hall's Theorem]

Let $G$ be a bipartite graph partited as $(X, Y)$. There exists a matching $M$ such that saturate all vertices in $X$ if and only of $\forall X' \subseteq X, |N(X')\cap Y| \geq |X'|$ holds. 


\end{tcolorbox}

\begin{tcolorbox}[colback=yellow!10!white,colframe=red!75!black,title=K\"onig's Theorem]

Vertex covering and Matching is dual in bipartite graphs. i.e., 
the size of optimal vertex cover and maximum matching is same on the bipartite graph. 

\end{tcolorbox}

From now on, we'll explore various properties on crown decomposition and use it to construct kernelization. 

\begin{tcolorbox}[colback=yellow!10!white,colframe=green!75!black,title=Lemma]

If $(R, H, C)$ is a crown decomposition of $G$, then there exists an optimal vertex cover of $G$ containing all vertices of $H$. 

\end{tcolorbox}

\begin{proof}
    
Suppose that $S$ be the any vertex cover of $G$, whose crown decomposition is $(R, H, C)$. Then, 
w.r.t. such decomposition, we can also decompose our vertex cover $S$ into $(S_R, S_H, S_C)$. 
Note that, $E(H \setminus S_H, C \setminus S_C) = \emptyset$. So, it means that, $N(H \setminus S_H) \cap C \subseteq S_C$. 
Then, $|N(H \setminus S_H) \cap C| \leq |S_C| $. By Hall's condition, 
\[|H \setminus S_H| \leq |N(H \setminus S_H) \cap C| \leq |S_C|\]
And it is clear that $S' = S - S_C + (H \setminus S_H)$ is a vertex cover. Since 
\[|S'| = |S| - |S_C| + |H \setminus S_H| \leq |S|\]
holds, it implies that there exists an optimal vertex cover containing entire $H$. 

\end{proof}

Similar to previous chapter, this lemma gives us following reduction rule. 

\begin{tcolorbox}[colback=yellow!10!white,colframe=blue!75!black,title=Reduction Rule 1]

For input $(G, k)$, let the crown decomposition of $G$ be $(R, H, C)$. Then, output $(G - (H \cup C), k - |H|)$
i.e., 
\[(G, k) \mapsto (G - (H \cup C), k - |H|)\]

\end{tcolorbox}
Now, we can prove that this is safe easily. 

\begin{proof}
    
If $(G, k)$ is yes instance, then by above lemma, there exists optimal vertex cover $S$ containing entire $H$. 
Then, it is clear that $S_C = \emptyset$ and $S_R$ is optimal vertex cover of $G[R]$. So, $(G - (H \cup C), k - |H|) = (G[R], k - |H|)$ is also yes. 

Conversely, obvious. 

\end{proof}

Actually, you might realized that this reduction rules is same to previous section's one. 
What's the difference? Previous one is LP-based decomposition. As we've seen, it also generate crown decomposition. 
However, what we defined for crown decomposition is more general concept. So, one can expect is that 
crown decomposition algorithm is much faster than construct crown decomposition via solving half-integral LP. 
So we need to construct algorithm for general crown decomposition. 

To show how can we construct crown decomposition and kernelization by Reduction Rule 1 at once, 
first we propose the complete kernelization algorithm. (See next page)

\newpage 

\begin{algorithm}
    
\caption{$4k + 1$ Kernelization for \textsc{Vertex Cover} using crown decomposition}
\begin{algorithmic}[1]
\Procedure{VC-Crown}{$G, k$}
\State Remove all isolated vertices in $G$ 
\If{$|V(G)| \leq 4k + 1$} \Return $(G, k)$ \Comment{Kernelization Success} 
\Else
\State $M \gets$ Find a maximal matching via greedy \Comment{Just pick until can't}
\If{$|M| \geq k + 1$} \Return small No-instance $(G_0, k_0)$
\State $C_0 \gets V(G) - V(M)$,  $ H_0 \gets N(C_0)$, $R_0 = V(G) - (C_0 \cup H_0)$
\While{$\exists H' \subseteq H_i$ s.t. $|H'| > |N(H') \cap C_i|$} \Comment{(*)}
\State $C_{i+1} \gets C_i - (N(H') \cap C_i)$
\State $H_{i+1} \gets H_i - H'$
\State $R_{i+1} \gets R_i \cup H' \cup (N(H') \cap C_i)$
\State $i \gets i + 1$
\EndWhile
\State \Return \textsc{VC-Crown}($G-(H_i \cup C_i), k - |H_i|$)
\EndIf
\EndIf
\EndProcedure
\end{algorithmic}
\end{algorithm}
Here, suppose that (*) works well, i.e., it terminates in polynomial time and always return crown decomposition of $G$. 
Then, the safeness of this kernelization is proved. Here, what we need to show to claim this is valid kernelization is that 

\begin{tcolorbox}[colback=yellow!10!white,colframe=red!75!black,title=Crown Decomposition Algorithm]

    \begin{itemize}
        \item the result of while loop is crown decomposition,
        \item checking of while condition (*) can be done in polynomial time.
    \end{itemize}
    
\end{tcolorbox}

\begin{proof}
    
Let's see the first one. Here, we claim loop invariants first; 
\begin{enumerate}
    \item $H_i \cup C_i \neq \emptyset$
    \item $C_i$ is independent set of $G$
    \item $N(C_i) \subseteq H_i$, i.e., $H_i$ seperates $C_i$ and $R_i$
    \item $|H_i| < |C_i|$
\end{enumerate}

\begin{itemize}
    \item (Initial Case)
    \begin{enumerate}
        \item $H_0 = N(C_0)$. So, if $H_0 \cup C_0 = \emptyset$, it implies that $C = \emptyset$. It means that 
        $M$ is perfect matching of $G$. However, $|V(G)| > 4k + 2$, so if the perfect matching exists, $|M| > 2k + 1$, contradiction. So, $H_0 \cup C_0 \neq \emptyset$. 
        \item Since $M$ is maximal matching, $C_0 = V(G) - V(M)$ must be independent. 
        \item Since $H_0 = N(C_0)$, this is obvious. 
        \item We know that $|V(G)| > 4k + 2$ and $|M| \leq k$, so $|V(M)| \leq 2k$, so $|C_0| > 2k + 2$. However, 
        $H_0 = N(C_0) \subseteq V(M)$ since if not, contradiction to $M$ is maximal. So, $|H_0| \leq |V(M)| \leq 2k$. 
        It implies that $|H_0| < |C_0|$ holds. 
    \end{enumerate}
    \item (Induction Step)
    It is recommended : Draw how we derive next iteration. 
    \begin{enumerate}
        \item Since $|H_i| < |C_i|$, the entire $H_i$ never selected to $H'$. So, $H' \subsetneq H_i$, so $H_{i+1} \cup C_{i+1} \neq \emptyset$. 
        \item $C_{i+1} \subset C_i$, so $C_{i+1}$ is obviously independent set. 
        \item It is sufficient to show that $C_{i+1}$ and $H' \cup (N(H') \cap C_i)$ is independent. 
        It is clear that $C_{i+1}$ and $N(H') \cap C_i$ is independent, since $C_i$ is independent. And, $N(H') \cap C_{i+1} = \emptyset$ since we removed all neighbors of $H'$ in $C_i$. 
        \item $|H_i| < |C_i|$. For $|H_{i+1}|$, we removed $|H'|$ vertices. For $|C_{i+1}|$, we removed $|N(H') \cap C_i|$ vertices. Since $|H'| > |N(H') \cap C_i|$, automatically $|H_{i+1}| < |C_{i+1}|$ holds. 
    \end{enumerate}
\end{itemize}

It implies that, if this while loop terminate, (1)-(3) holds and no such $H'$ exists, so it becomes crown decomposition of $G$. 

Now, let's see the second one. Which states we can decide such $H'$ (if exists) in polynomial time. Here, 
there is a polynomial time algorithm for \textsc{Minimum Vertex Cover} problem on bipartite graph. 
Let's imagine bipartite graph $(H \cup C, E(H, C))$. And apply above algorithm here, 
then we can find minimum vertex cover $X$. 

If $H_i \subseteq X$, then implies $H_i = X$. Then there exists matching $M^*$ such that saturate all vertices in $H_i$. So, Hall's condition holds. 
Then we can't find $H'$. 

If $H_i \setminus X \neq \emptyset$, then let $H' = H_i \setminus X$. 
Here, $N(H') \subseteq X \cap C_i$. If $|X \cap C_i| > |H'|$, then we can find minimum vertex cover by alternating 
all vertices in $X \cap C_i$ into $H'$. So, $|H'| \leq |X \cap C_i|$. If same, we can imagine that $H_i = X$ case. If else, our choice is good. 

All these processes can be done in polynomial time without brute forces. 

\end{proof}


\newpage 

\section{$\CO(k^2)$-Vertex Kernel for \textsc{Feedback Vertex Set}}

This section require quite long proofs and theorems. So, I'll write the full algorithm first. 
Note that, in this problem we assume that $G$ \textbf{allows multigraphs even loops.}

\begin{algorithm}
    
\caption{$\CO(k^2)$-kernelization for Feedback Vertex Set}

\begin{algorithmic}[1]
\Procedure{FVS-Kernel}{$G, k$}
\If{$k < 0$} \Return $(G_0, k_0)$ \Comment{Tiny No instance}
\EndIf
\State $(G, k) \gets $ Closure under Reduction Rule 1-4.
\If{$|V(G)| \leq 7k^2 + 7k$} \Return $(G, k)$ \Comment{Kernelization Success}
\ElsIf{$\Delta(G) < 7k$} \Return $(G_0, k_0)$ \Comment{Tiny No instance}
\Else  
\State $x \gets x \in V(G), d_G(x) \geq 7k$
\If{$x$ has $k+1$ size $x$-flower} \Return $(G - x, k - 1)$ \Comment{Reduction Rule 5} 
\Else    
\State \textbf{Assert} : $\exists S_x \subseteq V(G) - x$, hitting all $x$-cycles, $|S_x| \leq 2k$
\State $\CC \gets $ all components of $G - S_x - x$, connected with $x$
\If{$\CC$ has more than $k$ cycles} \Return $(G_0, k_0)$
\Else   
\State $\CT \gets$ Tree components in $\CC$ joined with $x$ by unique edge.
\State $S', T' \gets$ apply 2-expansion Lemma with $S_x, \CT$
\State \Return $(G, k) \gets$ Apply reduction rule 6 with $S', T'$
\EndIf
\EndIf 
\EndIf
\EndProcedure
\end{algorithmic}

\end{algorithm}

\begin{tcolorbox}[colback=yellow!10!white,colframe=blue!75!white,title=Reduction Rules for FVS kernelization,breakable]

\begin{enumerate}
    \item If $x \in V(G)$ has loops, then $(G - x, k - 1)$
    \item If $x \in V(G), d_G(x) = 0, 1$, then $(G - x, k)$
    \item If $u, v \in V(G)$ is connected by more than 2 parallel edges, remaining only 2.
    \item If $x \in V(G), d_G(x) = 2, ux, vx \in E(G)$ then $(G - x + uv, k)$
    \item If $x \in V(G)$, $x$ has $k+1$-star, then $(G - x, k - 1)$
    \item Suppose that there exists $x \in V(G)$, $S \subseteq V(G) - x$, $\CT$ is subset of components in $G - S - x$ such that 
    \begin{itemize}
        \item $\forall \CC \in \CT, G[\CC]$ is tree.
        \item $\forall \CC \in \CT$, $x$ and $\CC$ is connected by unique edge. 
        \item $\forall Z \subseteq S$, at least $2|Z|$ components in $\CT$ is connected with $Z$.
    \end{itemize}
    Then, we obtain new kernel $(G', k')$ by 
    \begin{itemize}
        \item Remove all unique edges between $x$ and $\CT$. 
        \item Make $x$ and all vertices in $S$ connected via double edges. 
        \item $k' = k$
    \end{itemize}
\end{enumerate}

\end{tcolorbox}

Yes, this needs quite long context of proof. Let's do this step-by-step. 

\begin{lem}
    
Reduction Rules 1-4 is safe and we can find irreducible kernel $(G, k)$ in polynomial time. 
Then, $\delta(G) \geq 3$ holds.

\end{lem}

\begin{lem}
    
Suppose that $(G, k)$ be the above kernel(RR 1-4). If $(G, k)$ is Yes-instance of \textsc{Feedback Vertex Set}, 
then following always holds: 
\[\Delta(G) < d \implies |V| < (d+1)k, \quad |E| < 2dk\]
In other words, if it is vanished, $(G, k)$ is always No-instance.  

\end{lem}

\begin{proof}
    
Since $(G, k)$ is Yes instance and kernel of RR 1-4, we can obtain 
\[\exists X \subseteq V(G), \quad |X| \leq k, \quad G - X \text{ is forest }, \quad 3 \leq d_G(v) < d\]
Then, 
\begin{align*}
    3 |V - X| \leq \sum_{u \in G - X} d_G(u) &= \big| E(X, G-X) \big| + 2 \big| E(G-X) \big| \\
    & \leq dk + 2(|V - X| - 1) \tag{Remark 1} \\
    & \leq 2|V| + (d-2)k - 2 \\
    & < 2|V| + (d-2)k
\end{align*}

This implies that 
\[3|V| - 3k < 2|V| + (d-2)k \implies |V| < (d+1)k\]
For the edges, 
\begin{align*}
    |E| &= \text{edges incident with $X$} + \text{edges does not touch $X$} \\
    & \leq dk + |V - X| - 1 \\
    & < 2dk \tag{By Remark 1}
\end{align*}

\end{proof}

With these, we can run our algorithm until line 7. If size is bounded, then return that as a kernel.
If size is unbounded by $d = 7k$ into above lemma, it vanishes Lemma 3.4.2. directly. So $(G, k)$ is automatically No instance, 
so output tiny No-instance here. Now we arrive at line 7. Let's do further analysis with added condition: 
\[3 \leq \delta(G), 7k \leq \Delta(G), |V| > (7k + 1)k\]

\begin{tcolorbox}[colback=yellow!10!white,colframe=brown!75!black,title=$x$-cycle and $x$-flower]

\textit{$x$-cycle} is a cycle such that containing $x$. \textit{$x$-flower} is collection of $x$-cycles such that vertex disjoint 
except only $x$. The \textit{order} of an $x$-cycle is number of $x$-cycles in that. 

\end{tcolorbox}

With this notion, it is easy to show that line 8 is safe and run in polynomial time with respect to $n$. 

\begin{lem}
    
Reduction Rule 5 is safe and run in polynomial time.

\end{lem}

Then now, let's move on to line 10. Here we have that $x$ does not have $k+1$ size flower and meet an assertation: 
\[\exists S_x \subseteq V(G) - x, \text{ hitting all $x$-cycles, $|S_x| \leq 2k$}\]
We need to prove that this can be asserted always. 

\begin{tcolorbox}[colback=yellow!10!white,colframe=red!75!black,title=Want To Show]

We want to show that, at line 7, one of following must hold: 
\begin{itemize}
    \item $x$ has $k+1$ size of $x$-flower. 
    \item No more than $k$ flower, then $\exists S_x \subseteq V(G) - x$, hitting all $x$-cycles, $|S_x| \leq 2k$
\end{itemize}

\end{tcolorbox}

\begin{tcolorbox}[colback=yellow!10!white,colframe=brown!75!black,title=A-path]

For $A \subseteq V(G)$, we say that a path $P$ is $A$-path if $|P| \leq 1$ and only endpoints of $P$ are contained in $A$. 

\end{tcolorbox}

We can divide the neighbors of $x$ into 2 groups: 
\begin{itemize}
    \item $A_1$: Neighbors of $x$, connected with $x$ by one edge. 
    \item $A_2$: Neighbors of $x$, connected with $x$ by two edges. 
\end{itemize}
In this setting and notion of $A$-paths, we can construct following lemma. 

\begin{lem}
Let $(G, k)$ be any graph and $x \in V(G)$. Following holds.
\begin{itemize}
    \item Maximum order of $x$-flower is summation of $|A_2|$ and maximum number of vertex disjoint $A_1$ paths in $G - x - A_2$. 
    \item $S_x \subseteq V(G) - x$ hit all $x$-cycles if and only if $S_x$ contains $A_2$ entirely, and $S_x - A_2$ hit all $A_1$ paths in $G - x - A_2$. 
\end{itemize}

\end{lem}

\begin{proof}
    
Easy

\end{proof}

Since we can find $A_2$ directly, our problem is reduced as following. 

\begin{tcolorbox}[colback=yellow!10!white,colframe=red!75!black,title=Want To Show - Reduced Version]

    Let $l := k - |A_2|$ , $A := A_1$, $H = G - x - A_2$. 
    \begin{itemize}
        \item There exists $l + 1$ vertex disjoint $A$-paths in $H$. 
        \item No more than $l$ vertex disjoint $A$-paths, then $\exists S_x \subseteq V(H), S_x$ hit all the $A$-paths in $H$, and $|S_x| \leq 2l$.
    \end{itemize}
    
\end{tcolorbox}

We'll take Gallai's theorem in Graph Theory. 

\begin{tcolorbox}[colback=yellow!10!white,colframe=blue!75!white,title=Gallai's A-path Theorem]

Let $G = (V, E)$ and $A \subseteq V$. 
Let's denote the maximum number of vertex disjoint $A$-path in $G$ by $M_A$. 
\[M_A = \min_{X \subseteq V} \left( |X| + \sum_{C \in \CC(G-X)}\lfloor \frac{A \cap C}{2} \rfloor \right)\]
$\CC(G-X)$ is collection of all components in $G - X$. And this theorem is algorithmic. i.e., we can find 
the solution $X$ in polynomial time. 
\end{tcolorbox}

With this theorem, we can prove our goal in red box. 

\begin{lem}
    
Let $G = (V, E)$ be a graph. If there exists at most $l$ vertex disjoint $A$-paths, 
then we can find $Z$ such that hit all the $A$-paths and $|Z| \leq 2l$. 

\end{lem}

\begin{proof}
    
According to the Gallai's Theorem, we can find $X \subseteq V(G)$ such that 
\[|X| + \sum_{C \in \CC(G-X)} \lfloor \frac{|A\cap C|}{2} \rfloor \leq l\]

Now, we'll construct $Z$ from here. First, add all vertices in $X$ into $Z$. And then, 
for every component $C$ in $\CC(G - X)$, choose all $A$-vertices except at most one. Then such cardinality of chosen vertices is same with 
$2 \cdot \lfloor \frac{|A\cap C|}{2}\rfloor$. And then add them into $Z$. Then, 
\[|Z| = |X| + 2 \sum_{C \in \CC(G - X)} \lfloor \frac{|A\cap C|}{2} \rfloor \leq 2\left(|X| + \sum_{C \in \CC(G - X)} \lfloor \frac{|A\cap C|}{2} \rfloor\right) \leq 2l\]
Here, this $Z$ hits all $A$-paths. Since we choose quite amount of $A$-vertices, sufficient to consider remaining $A$-vertices and their paths. 
Remaining vertices are located in each components, at most one in there. So, if there exists any $A$-path, 
such path must through $X$ and reach to another component. Since $X \subseteq Z$, it is hit by $Z$. 
\end{proof}

\newpage

Now we arrive at line 11 safely. 
The next step will be reduction of given graph with $S_x$ where $S_x \subseteq V(G) - x$ hitting all $x$-cycles and $|S_x| \leq 2k$. 
There remains most treaky part here. Let's rewrite Reduction Rule 6 again. 

\begin{tcolorbox}[colback=yellow!10!white,colframe=blue!75!black,title=Reduction Rule 6]

Apply condition: there exists $x \in V(G), S \subseteq V - x$ and $\CT \subseteq \CC(G - x - S)$ s.t.
\begin{itemize}
    \item $G[C]$ is tree, for each $C \in \CT$
    \item There exists unique edge between $x$ and each $C$
    \item $\forall Z \subseteq S$, number of incident components in $\CT$ with $Z$ is more than $2|Z|$ 
\end{itemize}

Then, reduce the instance $G$ by: 

\begin{itemize}
    \item Remove all edges between $x$ and $\CT$
    \item Make every vertices in $S$ is adjacent with $x$ via exactly 2 edges.
\end{itemize}

\end{tcolorbox}

\begin{lem}
    
Reduction Rule 6 is safe.

\end{lem}

\begin{proof}
    since we changes only the edges incident with $x$, only $x$-cycles are matter. 

($\Rightarrow$) Suppose that $(G, k)$ is Yes instance of FVS. Here, let $X$ be an optimal 
FVS of $G$. If $x \in X$, since $x$ hit all the newly generated $x$-cycles, so $X$ is also a FVS of $G'$. 
If $x \notin X$, we'll find another optimal FVS $X^*$ such that $S \subset X$. If we can find, then 
since $S$ hit all the newly generated $x$-cycles in $G'$, $X^*$ become FVS of $G'$. 
Let $Z = S \setminus X \neq \emptyset$. For each $z \in Z$, at least 2 components ($\CT_1, \CT_2$) are incident with $z$ in $\CT$ (Apply Condition 3)
Then since we can find $z \rightarrow \CT_1 \rightarrow x \rightarrow \CT_2$ cycle, $X$ must hit this in $\CT_1$ or $\CT_2$.
Let's remove all the vertices of $X$ in $\CT$ and add all vertices in $Z$. 
\[X^* = \left(X \setminus \bigcup_{T \in \CT} V(T) \right) \cup Z\]
Here, since amount of neighbor components of $Z$ is larger than $2|Z|$, at least $|Z|$ vertices of $|X|$ are included in $\CT$. Since 
we remove all them and add $Z$, $|X^*| \leq |X|$. And $X^*$ is clearly FVS since all cycles s.t. vertices in $\CT$ covered were hit also vertices in $Z$. So 
$Z$ hit all the cycles that $X \setminus \bigcup_{T \in \CT} V(T)$ covered, so $X^*$ is FVS. 
Then $X^*$ is also FVS of $G'$. 

($\Leftarrow$) Here, for any FVS $X'$ of $G'$, $x \in X'$ or $S \subseteq X'$. So both case, optimal FVS $X'$ of 
$G'$ also be the FVS of $G$. 

\end{proof}

\newpage 

One we must show is that, we can apply Reduction Rule with conditions until line 14 always. 
Actually, it is possible. 

\begin{lem}
    
Let $H$ be a bipartite graph $S \uplus T$. In polynomial time, one can find a set $Z \subset S$ with 
$|N(Z)| < 2|Z|$ if such $Z$ exists. 

\end{lem}

\begin{proof}
Using vertex copy of $S$ then K\"onig's theorem and obtain maximum matching. 
\end{proof}

\begin{lem}
    
Let $H$ be a bipartite graph with a bipartition $S \uplus T$ s.t. 
\begin{itemize}
    \item every vertex of $T$ adjacent with vertex of $S$
    \item $|T| \geq 2|S|$
\end{itemize}
Then, in polynomial time one can find vertex sets $S' \subseteq S$ and $T' \subseteq T$ such that 
\begin{itemize}
    \item $N(T') = S'$ and 
    \item $|N(Z)| \geq 2|Z|$ for every $Z \subseteq S'$
\end{itemize}

\end{lem}

\begin{lem}
    Line 14-16 is always applicable.
\end{lem}

\begin{proof}
    
We have $S_x \subseteq V(G) - x$ such that $|S_x| \leq 2k$ and $S_x$ hits all $x$-cycles. 
However, $d_G(x) \geq 7k$. Since $S_x$ contains all double edges neighbors of $x$, but there limit is $k$ since no order $k+1$ $x$-flower condition. Implies that, at least $4k$ are not in $S_x$. Moreover, since they have not to construct any $x$-cycles, 
each of them are incident with tree components in $G - x - S_x$ via unique edges. Here, by imagine that such tree components are vertices, 
and $S$, they forms bipartite graph. Note that $|\CT| \geq 2|S_x|$ and every $T \in \CT$ is adjacent with at least one vertex of $S$, because 
in $G$, $T$ does not have leaf. So we can apply Lemma 3.4.8, and obtain $S' \subseteq S_x$, $\CT' \subseteq \CT$. 
They holds the Apply Condition of Reduction Rule 6.

\end{proof}

\chapter{Iterative Compression}

\section{$\CO(5^k)$-time algorithm for Feedback Vertex Set}

\begin{tcolorbox}[colback=yellow!10!white,colframe=brown!75!black,title=Compression FVS]

\textsc{Compression FVS}

\textbf{Input}: Graph $G = (V, E)$ and Feedback Vertex Set $X$ of $G$.

\textbf{Parameter}: $|X|$

\textbf{Output(Decision)}: Does $G$ have another FVS $Y$ such that $|Y| < |X|$?

\end{tcolorbox}

The Iterative Compression method is that, solve original problem via solving its compression version. 
In this case, we'll study Compression version of \textsc{Feedback Vertex Set} problem.
Suppose that our graph is evolving by one vertex.
\[G_1 \rightarrow G_2 \rightarrow \cdots \rightarrow G_n = G\]
If we know $X_i$, the Feedback Vertex Set of $G_i$, then obviously $X_{i+1} = X_i \cup \{v_{i+1}\}$ become 
feedback vertex set of $G_{i+1}$. If it doesn't exceed our original budget, it can be the instance for $G_{i+1}$. 
However, if it exceed our budget, we need to solve \textsc{Compression FVS}$(G_{i+1}, X_{i+1})$. 
This is sort of dynamic programming, let's formalize this method based on FVS problem. 

\begin{lem}

If we can solve \textsc{Compression FVS}$(G, X)$ in time $c^{|X|} \cdot n^d$, then we can solve 
\textsc{Feedback Vertex Set}$(G, k)$ in time $c^k \cdot n^{d+1}$.

\end{lem}

\begin{proof}
    
Let's enumerate vertices of $G$ by 
\[v_1, v_2, \cdots, v_n, \quad G_i := G[\{v_1, \cdots, v_i\}]\]
We construct algorithm that we mentioned above. The worst case running time of \textsc{Compression FVS}$(G, X)$ is 
$c^k \cdot n^d$ since $|X| \leq k$ in our setting. So, the worst case running time of our algorithm is that 
\[n \times (c^k\cdot n^d) = c^k \cdot n^{d+1}\]

\end{proof}

Thanks to this lemma, we are sufficient to show that we can solve \textsc{Compression FVS} in such time complexity. 
In similar way, we'll reduce \textsc{Compression FVS} again into another easier problem. 

\begin{tcolorbox}[colback=yellow!10!white,colframe=brown!75!black,title=Disjoint FVS]

\textsc{Disjoint FVS}

\textbf{Input}: Graph $G = (V, E)$, $X$ is a feedback vertex set of $G$, $k \in \BZ$

\textbf{Parameter}: $k + |\CC(G[X])|$

\textbf{Output}: Does $G$ have disjoint FVS $Y$ (i.e., $X \cap Y = \emptyset$) s.t. $|Y| \leq k$?

\end{tcolorbox}

\begin{algorithm} 
\caption{Algorithm for Disjoint FVS}
\begin{algorithmic}[1]
\Procedure{Disjoint FVS}{$G, X, k$}
\State $G \gets$ delete all $d_G(v) = 0,1$ vertices  
\State $F \gets G[V \setminus X]$, $\CC \gets$ components of $G[X]$ \Comment{$F$ is forest}
\State $\forall x \in F, d_G(x) = 2, |N(x) \cap F| = 1$ then bypass $x$.
\If{$G$ is acyclic} \Return Yes, by witness : $\emptyset$
\EndIf
\If{$G[X]$ contains any cycle} \Return No      
\EndIf
\If{$G$ has any cycle and $k <= 0$} \Return No
\EndIf
\State $l \gets $ any leaf vertex in $F$
\If{$\exists C \in \CC$, $l$ incident $C$ with more than 2 edges}
\State \Return \textsc{Disjoint FVS}$(G - l, X, k-1) \cup \{l\}$
\ElsIf{No such cases}
\State \Return \textsc{Disjoint FVS}$(G - l, X, k-1) \cup \{l\}$ $\vee$ \textsc{Disjoint FVS}$(G, X \cup \{l\}, k)$
\EndIf
\EndProcedure
\end{algorithmic}

\end{algorithm}

\begin{lem}
    
This algorithm is correctly solves \textsc{Disjoint FVS} in time $\CO^*(2^{\mu(I)})$ where $\mu(I) = k + |\CC(G[X])|$

\end{lem}

\begin{proof}
    
Since our goal is find feedback vertex set disjoint with $X$, 
line 5,6,7 is clearly correct. And for reduction step, if such $x$ in line 4 exists, then 
when we need to pick $x$, then instead pick its $F$-neighbor, so this also clear. 
Then suppose that $l$ is leaf in $F$. If line 9 occurs, then $l$ joined in a cycle with such $\CC$. So we must pick $l$ as new FVS. 
If not occurs, we must consider 2 cases each: $Y$ contains $l$ or not. This is line 12.

For the running time analysis, 
\[\mu(I = (G, X, k)) = k + |\CC(G[X])|\]
For each step, if we choose $I' := (G - l, X, k - 1)$ then $\mu(I') = \mu(I) - 1$ so decreases, 
if we choose $I' := (G, X \cup \{l\}, k)$, since in line 12, $l$ works as \lq bridge' of components of $G[X]$ so 
$\mu(I') \leq \mu(I) - 1$. So our parameter always strictly decreases at least 1, and the worst branch cases are 2, so the running time would be $\CO^*(2^{\mu(I)})$. 
This completes proof.

\end{proof}

\begin{lem}
    
With above algorithms, we can solve \textsc{Compression FVS} in time $5^{|X|}\cdot n^d$

\end{lem}

\begin{algorithm}
\caption{Algorithm for Compression FVS}
\begin{algorithmic}[1]
\Procedure{Compression FVS}{$G, X$}
\For{$\forall I \subsetneq X$} 
\State $\tilde{G} \gets G - I, \tilde{X} \gets X \setminus I, \tilde{k} \gets |\tilde{X}| - 1$
\If{\textsc{Disjoint FVS}$(\tilde{G}, \tilde{X}, \tilde{k})$ returns No} Continue
\Else 
\State $\tilde{Y} \gets \textsc{Disjoint FVS}(\tilde{G}, \tilde{X}, \tilde{k})$
\State \Return Yes with witness: $I \cup \tilde{Y}$
\EndIf
\EndFor
\State \Return No
\EndProcedure
\end{algorithmic}
\end{algorithm}

\begin{proof}
    
It is obvious that, if this algorithm returns Yes, then such instance $I \cup \tilde{Y}$ is actually 
smaller size FVS of $G$. This is because we fix $I$, and then find alternatives for $X \setminus I$ vertices to hit all cycles that 
$X \setminus I$ hit. So this algorithm is complete. How about soundness? Suppose that this algorithm returns No. 
If the instance $(G, X)$ is actually Yes-instance, then there exists smaller FVS $Y \neq X$ of $G$. 
Then, we can find $Y \cap X$. Since $I$ iterate all proper subsets of $X$, $I = X \cap Y$ is iterated. 
Since it doesn't return Yes, so \textsc{Disjoint FVS}$(G - X\cap Y, X - Y, |X - Y| - 1)$ returns No. 
However, we have remaining parts of $Y$, which is $Y - X$ such that $|Y - X| < |X - Y|$. Since $Y - X$ 
is FVS of $G - X \cap Y$, this is contradiction to \textsc{Disjoint FVS} is correct. So above algorithm is correct. 

Finally, let's analyze the running time of this algorithm. For each iteration of $I$, 
\[\mu(G - I, X \setminus I, |X \setminus I| - 1) = |X \setminus I| - 1 + |\CC(\tilde{G}[\tilde{X}])| \leq 2(|X| - i) - 1\]
So, for the selected subset size $i$, running time will be $\CO^*(2^{2(|X|-i)-1}) = \CO^*(4^{|X|-i})$
So, running time will be 
\[\sum_{i=0}^{|X|-1}\binom{|X|}{i} \cdot \CO^*(4^{|X|-i}) \leq (4 + 1)^{|X|}\cdot n^d\]


\end{proof}

\chapter{Randomized Algorithms}

\chapter{Tree Decomposition}

\begin{tcolorbox}[colback=yellow!10!white,colframe=brown!75!black,title=Tree Decomposition]
Let $G$ be a graph. A \textit{tree decomposition} of $G$ is pair $(T, \chi)$ where $T$ is a tree 
and $\chi: V(T) \rightarrow 2^{V(G)}$. For each $t \in V(T)$, we say $\chi(t) \subseteq V(G)$ be a \textit{bag} and $(T, \chi)$ satisfy following. 
\begin{enumerate}
    \item $\bigcup_{t \in V(T)} \chi(t) = V(G)$
    \item $\forall e = uv \in E(G), \quad \exists ~ t \in V(T), ~ \{u,v\} \subseteq \chi(t)$
    \item $\forall v \in V(G)$, $\{t \in V(T)~|~ v \in \chi(t)\}$ forms a subtree of $T$. 
\end{enumerate}
\end{tcolorbox}

Note that, this definition is somewhat similar to \textit{topology} on sets. We define the measure that how finer our tree decomposition is, called \textit{width}. 

\begin{defn}
For each tree decomposition $(T, \chi)$ of $G$, we define \textit{width} by $\max_{t \in V(T)} |\chi(t)| - 1$. 
i.e., the largest bag size - 1. And we denote the global measure of $G$ by \textit{treewidth} of $G$, defined by 
\[tw(G) = \min_{(T, \chi)} \text{width}(T, \chi)\]
\end{defn}

\begin{defn}
Suppose that $A, B \subseteq V(G)$. We define $(A,B)$-path in $G$ as a path with one endpoint is in $A$ and another is in $B$. 
Note that $v \in A \cap B$ is trivial $(A, B)$-path. We call that $X$ is an $(A,B)$-seperator in $G$ if 
every $(A, B)$-path contains a vertex of $X$. We define \textit{contraction} of $G$ by $e = uv \in E(G)$ is an graph modification such that 
merge $u, v$. Formally, delete $u, v$ and add new vertex $z_{uv}$ and then connect with $N_G(u) \cup N_G(v)$. Denote such modified graph by $G/e$. 
\end{defn}

\begin{tcolorbox}[colback=yellow!10!white,colframe=blue!75!black,title=Propositions,breakable]
    \begin{enumerate}
        \item[(1)] $tw(G) - 1 \leq tw(G - v) \leq tw(G)$
        \item[(2)] $tw(G) - 1 \leq tw(G - e) \leq tw(G)$
        \item[(3)] $tw(G) - 1 \leq tw(G/e) \leq tw(G)$
        \item[(4)] $G$ is a forest $\iff$ $tw(G) \leq 1$  
    \end{enumerate}
\end{tcolorbox}

\begin{tcolorbox}[colback=yellow!10!white,colframe=green!75!black,title=Nice Tree Decomposition]

A nice tree decomposition $(T, \chi)$ of $G$ is a rooted tree decomposition such that each $t \in V(T)$ is one of the following types: 
\begin{itemize}
    \item Leaf node: $t$ is a leaf node of $T$
    \item Introduce node: $t$ has exactly one child $t'$ holds that $\chi(t) = \chi(t') \cup \{v\}$ for some $v \in V(G)$. 
    \item Forget node: $t$ has exactly one child $t'$ and holds that $\chi(t) = \chi(t') \setminus \{v\}$ for some $v \in V(G)$. 
    \item Join node: $t$ has exactly 2 children $t_1, t_2$ such that $\chi(t) = \chi(t_1) = \chi(t_2)$. 
\end{itemize}

\end{tcolorbox}

\begin{lem}
    Let $(T, \chi)$ be a tree decomposition of $G$ of width $w$. In linear time, we transform this into a nice tree decomposition $(T', \chi')$ of $G$ of width $w$. Moreover, 
    \begin{itemize}
        \item Root node and Leaf nodes are empty bags. (If we want)
        \item $|V(T')| \leq 4wn$. 
    \end{itemize}
\end{lem}

\begin{proof}
    Let root be empty bag and for every leaf node of $T$ we append empty bag. Choose any $t \in V(T)$. 
    Following that, if $t$ has $d$ children then we construct $d$-times of join nodes and connect each children into unique child of generated join nodes. 
    If $t$ has unique child, remove $\chi(t) \setminus \chi(t')$ one-by-one via introduce nodes and add $\chi(t') \setminus \chi(t)$ one-by-one via forget nodes. 
\end{proof}

\begin{defn}
    We define that $T_t$ be the subtree of $T$ rooted at $t$. 
    \[V_t = \bigcup_{b \in V(T_t)} \chi(b) \quad G_t = G[V_t]\]
\end{defn}

\begin{tcolorbox}[colback=yellow!10!white,colframe=blue!75!black,title=Lemma]

Let $(T, \chi)$ be a nice tree decomposition of $G$. Suppose that $t \in V(T)$ is an 
introduce node with its unique child $t'$ such that $\chi(t) = \chi(t') \cup \{v\}$ for some $v \in V(G)$. 
Then, 
\[N_{G_t}(v) \subseteq \chi(t)\]

\end{tcolorbox}

\begin{proof}
    $T_t$ become a tree decomposition of $G_t$. However, the root is unique node of tree decomposition 
    such that containing $v$, so every edges incident with $v$ in $G_t$ is contained in $t$. This 
    exactly implies $N_{G_t}(v) \subseteq \chi(t)$
\end{proof}

\begin{tcolorbox}[colback=yellow!10!white,colframe=blue!75!black,title=Lemma]

Let $(T, \chi)$ be a nice tree decomposition of $G$. Suppose that $t \in V(T)$ is an 
join node with $t_1, t_2$ then $V_{t_1} \cap V_{t_2} = \chi(t)$. 

\end{tcolorbox}

\begin{proof}
    Note that vertices appears in $V_{t_1} \cap V_{t_2}$ must be contained in $\chi(t)$ since 
    $(T, \chi)$ is a tree decomposition. 
\end{proof}

\begin{tcolorbox}[colback=yellow!10!white,colframe=red!75!black,title=Maximum Independent Set]

For a graph $G = (V, E)$, find a maximum size of independent set of $G$. 

\end{tcolorbox}

\begin{thm}
There is an algorithm which, given $G$ and non-redundant tree decomposition $(T, \chi)$ such that width $\leq w$, 
finds a maximum size of independent set of $G$ in time $\CO(w2^w \cdot n)$. 
\end{thm}

\begin{tcolorbox}[colback=yellow!10!white,colframe=gray!75!black,title=Step 0]

We first transform $(T, \chi)$ into a \textit{nice tree decomposition} via Lemma 6.3. 

\end{tcolorbox}

\begin{tcolorbox}[colback=yellow!10!white,colframe=gray!75!black,title=Step 1]

Define following for $t \in V(T)$ and $S \subseteq \chi(t)$: 
\[\textsc{IND}_t[S] = \begin{cases}
    \text{Size of maximum independent set $I$ of $G_t$ such that $I \cap \chi(t) = S$} \\
    \perp \text{ if no such $I$ exists }
\end{cases}\]
Remark that Maximum size of independent set of $G$ is $\textsc{IND}_r[\emptyset]$. 
\end{tcolorbox}

\begin{tcolorbox}[colback=yellow!10!white,colframe=violet!75!white,title=Planning Mode]
    \[\text{(Initialization) At each leaf $t$ of $T$, compute a basis cases for DP.}\] 
    \[\Downarrow\]
    \[\text{(Recurrence) For forget/introduce/join node $t$, define recursive relation}\]
    \[\Downarrow\]
    \[\text{(Runtime efficiency) Compute $\textsc{IND}_t[S]$ for each $S \subseteq \chi(t)$ efficiently.}\]
\end{tcolorbox}

\begin{rem}
Note that our algorithm answer with non-witness i.e., does not compute actual maximum independent set. 
\end{rem}

\begin{tcolorbox}[colback=yellow!10!white,colframe=gray!75!black,title=Initialization]

Let $t$ be any leaf node of $T$. For every $S \subseteq \chi(t)$, we can easily compute 
\[\textsc{IND}_t[S] = \begin{cases}
    |S| & S \text{ is an independent set} \\
    \perp & \text{ otherwise }
\end{cases}\]

\end{tcolorbox}

\begin{tcolorbox}[colback=yellow!10!white,colframe=gray!75!black,title=Recurrence]

We construct following relations for each nodes' type. 

\end{tcolorbox}

If $t$ is a forget node, then the child $t'$ holds that $\chi(t) = \chi(t') \setminus \{v\}$. For $S \subseteq \chi(t)$, 
\[\]

\end{document}